% GENERATED FILE. Edit canonical lessons, then run npm run generate:books.
% canonical-source-hash: fnv1a64:94fc816bdc384821
% canonical-lessons: TE-C482-vadagallu, TE-C482-mancugadda, TE-C482-himapatam, TE-C482-pillagali, TE-C482-callabadu

\chapter{Cold Weather}
\label{ch:te-cold-weather}
\hlchaptermodality{482}

\begin{chapteropening}
\textbf{By the end of this chapter:} \emph{I can say hailstones, ice, snowfall, a breeze and to cool down.}

Complete the last lesson of chapter 482: I can say hailstones, ice, snowfall, a breeze and to cool down.
\end{chapteropening}

\section[\texorpdfstring{\te{వడగళ్ళు} (vaḍagaḷḷu)}{వడగళ్ళు (vaḍagaḷḷu)}]{\te{వడగళ్ళు} (vaḍagaḷḷu) — hailstones}

\label{lesson:TE-C482-vadagallu}



\begin{quote}
Before the new one: say the Telugu for to dry up, then the Telugu for to crack open.
\end{quote}

\subsection*{You'll want to know: \te{వడగళ్ళు}}
\textbf{\te{వడగళ్ళు}} — \emph{vaḍagaḷḷu} — \textquotedblleft{}hailstones\textquotedblright{}.

\begin{quote}
\te{నిన్న} \te{వడగళ్ళ} \te{వాన} \te{పడింది}.
\end{quote}

\begin{quote}
\emph{ninna vaḍagaḷḷa vāna paḍindi.}
\end{quote}

\begin{quote}
Yesterday there was a hailstorm.
\end{quote}

Plural: \te{వడగళ్ళు}, hailstones; \te{వడగళ్ళ} \te{వాన}, a hailstorm. -\te{గళ్ళు} is the plural of \te{కల్లు}, a stone.

\begin{cousinweb}
\textbf{\te{వడగళ్ళు}} (\emph{vaḍagaḷḷu}) is a Telugu compound: -\te{గళ్ళు}, \textquotedblleft{}stones\textquotedblright{}, from the native \te{కల్లు}, \textquotedblleft{}a stone\textquotedblright{} (DEDR 1298), and \te{వడ}-, whose origin is not settled.
\end{cousinweb}

\begin{grammarlens}[title={What you've built}]
One more word for getting around.
\end{grammarlens}

\subsection*{Your turn}
\begin{itemize}
\raggedright
  \item \emph{Say it:} \emph{vaḍagaḷḷu}
  \item \emph{Say it:} \emph{vaḍagaḷḷu}, once more
  \item \emph{Say it:} \emph{ninna vaḍagaḷḷa vāna paḍindi.} — Yesterday there was a hailstorm.
  \item \emph{Recall:} say the Telugu for to dry up, then the Telugu for to crack open, then say \emph{vaḍagaḷḷu} again
\end{itemize}

\begin{tcolorbox}[breakable,colback=teal!4,colframe=teal!35!black,title={Before you move on}]
What does \te{వడగళ్ళు} mean? (\textquotedblleft{}Hailstones\textquotedblright{}.) Say it once more.
\end{tcolorbox}
\section[\texorpdfstring{\te{మంచుగడ్డ} (mañcugaḍḍa)}{మంచుగడ్డ (mañcugaḍḍa)}]{\te{మంచుగడ్డ} (mañcugaḍḍa) — ice; a block of ice}

\label{lesson:TE-C482-mancugadda}



\begin{quote}
Before the new one: say the Telugu for to crack open, then the Telugu for hailstones.
\end{quote}

\subsection*{You'll want to know: \te{మంచుగడ్డ}}
\textbf{\te{మంచుగడ్డ}} — \emph{mañcugaḍḍa} — \textquotedblleft{}ice; a block of ice\textquotedblright{}.

\begin{quote}
\te{చేపలను} \te{మంచుగడ్డల} \te{మధ్య} \te{పెట్టారు}.
\end{quote}

\begin{quote}
\emph{cēpalanu mañcugaḍḍala madhya peṭṭāru.}
\end{quote}

\begin{quote}
They packed the fish in ice.
\end{quote}

\te{మంచు}, snow, frost, + \te{గడ్డ}, a lump. Ice cubes for a drink are also \te{మంచుగడ్డలు}.

\begin{cousinweb}
\textbf{\te{మంచుగడ్డ}} (\emph{mañcugaḍḍa}) is native Dravidian: \te{మంచు}, \textquotedblleft{}mist, dew, frost\textquotedblright{} (DEDR 4641), and \te{గడ్డ}, \textquotedblleft{}a lump, a clod\textquotedblright{} (DEDR 1148).
\end{cousinweb}

\begin{grammarlens}[title={What you've built}]
One more word for getting around.
\end{grammarlens}

\subsection*{Your turn}
\begin{itemize}
\raggedright
  \item \emph{Say it:} \emph{mañcugaḍḍa}
  \item \emph{Say it:} \emph{mañcugaḍḍa}, once more
  \item \emph{Say it:} \emph{cēpalanu mañcugaḍḍala madhya peṭṭāru.} — They packed the fish in ice.
  \item \emph{Recall:} say the Telugu for to crack open, then the Telugu for hailstones, then say \emph{mañcugaḍḍa} again
\end{itemize}

\begin{tcolorbox}[breakable,colback=teal!4,colframe=teal!35!black,title={Before you move on}]
What does \te{మంచుగడ్డ} mean? (\textquotedblleft{}Ice; a block of ice\textquotedblright{}.) Say it once more.
\end{tcolorbox}
\section[\texorpdfstring{\te{హిమపాతం} (himapātaṁ)}{హిమపాతం (himapātaṁ)}]{\te{హిమపాతం} (himapātaṁ) — snowfall}

\label{lesson:TE-C482-himapatam}



\begin{quote}
Before the new one: say the Telugu for hailstones, then the Telugu for ice.
\end{quote}

\subsection*{You'll want to know: \te{హిమపాతం}}
\textbf{\te{హిమపాతం}} — \emph{himapātaṁ} — \textquotedblleft{}snowfall\textquotedblright{}.

\begin{quote}
\te{కొండల} \te{మీద} \te{భారీ} \te{హిమపాతం} \te{కురిసింది}.
\end{quote}

\begin{quote}
\emph{koṇḍala mīda bhārī himapātaṁ kurisindi.}
\end{quote}

\begin{quote}
Heavy snow fell on the mountains.
\end{quote}

\te{హిమ}, snow, + \te{పాతం}, a fall.

\begin{cousinweb}
\textbf{\te{హిమపాతం}} (\emph{himapātaṁ}) is a compound of two words borrowed from Sanskrit: \emph{hima}, \textquotedblleft{}snow, frost\textquotedblright{}, and \emph{pāta}, \textquotedblleft{}a fall\textquotedblright{} (Monier-Williams).
\end{cousinweb}

\begin{grammarlens}[title={What you've built}]
One more word for getting around.
\end{grammarlens}

\subsection*{Your turn}
\begin{itemize}
\raggedright
  \item \emph{Say it:} \emph{himapātaṁ}
  \item \emph{Say it:} \emph{himapātaṁ}, once more
  \item \emph{Say it:} \emph{koṇḍala mīda bhārī himapātaṁ kurisindi.} — Heavy snow fell on the mountains.
  \item \emph{Recall:} say the Telugu for hailstones, then the Telugu for ice, then say \emph{himapātaṁ} again
\end{itemize}

\begin{tcolorbox}[breakable,colback=teal!4,colframe=teal!35!black,title={Before you move on}]
What does \te{హిమపాతం} mean? (\textquotedblleft{}Snowfall\textquotedblright{}.) Say it once more.
\end{tcolorbox}
\section[\texorpdfstring{\te{పిల్లగాలి} (pillagāli)}{పిల్లగాలి (pillagāli)}]{\te{పిల్లగాలి} (pillagāli) — a breeze}

\label{lesson:TE-C482-pillagali}



\begin{quote}
Before the new one: say the Telugu for ice, then the Telugu for snowfall.
\end{quote}

\subsection*{You'll want to know: \te{పిల్లగాలి}}
\textbf{\te{పిల్లగాలి}} — \emph{pillagāli} — \textquotedblleft{}a breeze\textquotedblright{}.

\begin{quote}
\te{సాయంత్రం} \te{పిల్లగాలి} \te{హాయిగా} \te{ఉంది}.
\end{quote}

\begin{quote}
\emph{sāyantraṁ pillagāli hāyigā undi.}
\end{quote}

\begin{quote}
The evening breeze is pleasant.
\end{quote}

\te{పిల్ల}, little, young, + \te{గాలి}, wind: a gentle little wind.

\begin{cousinweb}
\textbf{\te{పిల్లగాలి}} (\emph{pillagāli}) is native Dravidian: \te{పిల్ల}, \textquotedblleft{}young, small\textquotedblright{} (DEDR 4198), and \te{గాలి}, \textquotedblleft{}wind\textquotedblright{} (DEDR 1499).
\end{cousinweb}

\begin{grammarlens}[title={What you've built}]
One more word for getting around.
\end{grammarlens}

\subsection*{Your turn}
\begin{itemize}
\raggedright
  \item \emph{Say it:} \emph{pillagāli}
  \item \emph{Say it:} \emph{pillagāli}, once more
  \item \emph{Say it:} \emph{sāyantraṁ pillagāli hāyigā undi.} — The evening breeze is pleasant.
  \item \emph{Recall:} say the Telugu for ice, then the Telugu for snowfall, then say \emph{pillagāli} again
\end{itemize}

\begin{tcolorbox}[breakable,colback=teal!4,colframe=teal!35!black,title={Before you move on}]
What does \te{పిల్లగాలి} mean? (\textquotedblleft{}A breeze\textquotedblright{}.) Say it once more.
\end{tcolorbox}
\section[\texorpdfstring{\te{చల్లబడు} (callabaḍu)}{చల్లబడు (callabaḍu)}]{\te{చల్లబడు} (callabaḍu) — to cool down}

\label{lesson:TE-C482-callabadu}



\begin{quote}
Before the new one: say the Telugu for snowfall, then the Telugu for a breeze.
\end{quote}

\subsection*{You'll want to know: \te{చల్లబడు}}
\textbf{\te{చల్లబడు}} — \emph{callabaḍu} — \textquotedblleft{}to cool down\textquotedblright{}.

\begin{quote}
\te{వర్షం} \te{తరువాత} \te{వాతావరణం} \te{చల్లబడింది}.
\end{quote}

\begin{quote}
\emph{varṣaṁ taruvāta vātāvaraṇaṁ callabaḍindi.}
\end{quote}

\begin{quote}
After the rain the weather cooled down.
\end{quote}

Past (I …): \te{చల్లబడ్డాను}, \emph{callabaḍḍānu}.

\te{చల్ల}, cool, + \te{పడు}. Tempers also cool this way.

\begin{cousinweb}
\textbf{\te{చల్లబడు}} (\emph{callabaḍu}) is native Dravidian: \te{చల్ల}, \textquotedblleft{}cool, coolness\textquotedblright{} (DEDR 2408), with \te{పడు}, \textquotedblleft{}to undergo\textquotedblright{} (DEDR 3853).
\end{cousinweb}

\begin{grammarlens}[title={What you've built}]
One more word for getting around.
\end{grammarlens}

\subsection*{Your turn}
\begin{itemize}
\raggedright
  \item \emph{Say it:} \emph{callabaḍu}
  \item \emph{Say it:} \emph{callabaḍu}, once more
  \item \emph{Say it:} \emph{varṣaṁ taruvāta vātāvaraṇaṁ callabaḍindi.} — After the rain the weather cooled down.
  \item \emph{Recall:} say the Telugu for snowfall, then the Telugu for a breeze, then say \emph{callabaḍu} again
\end{itemize}

\begin{tcolorbox}[breakable,colback=teal!4,colframe=teal!35!black,title={Before you move on}]
What does \te{చల్లబడు} mean? (\textquotedblleft{}To cool down\textquotedblright{}.) Say it once more.
\end{tcolorbox}
